%=======================================================================
%  Temporary Help Services in the United States
%  Beamer slides summarising the QWI analysis in this repository.
%
%  Build:  pdflatex temp_agencies_slides.tex   (run twice for the outline)
%  Figures are read from ../output/ (generated by the scripts in ../code/).
%=======================================================================
\documentclass[aspectratio=169,10pt]{beamer}

\usetheme{Boadilla}
\usecolortheme{seahorse}
\useinnertheme{rectangles}
\setbeamertemplate{navigation symbols}{}
\setbeamertemplate{footline}[frame number]
\setbeamertemplate{itemize items}[circle]
\setbeamercolor{block title}{bg=structure!25,fg=black}

\usepackage[T1]{fontenc}
\usepackage{lmodern}
\usepackage{booktabs}
\usepackage{siunitx}
\usepackage{graphicx}
\usepackage{caption}
\captionsetup{font=scriptsize,labelformat=empty}

\newcommand{\fig}[2][height=0.74\textheight]{%
  \begin{center}\includegraphics[#1]{../output/#2}\end{center}}
\newcommand{\source}[1]{{\scriptsize\textcolor{gray}{Source: #1}}}

\title[Temp Help Services (QWI)]{Temporary Help Services in the United States}
\subtitle{Employment, wages, and demographics from the Census Quarterly Workforce Indicators}
\author[QWI analysis]{Prepared from U.S.\ Census QWI microdata}
\date{\today}

\begin{document}

\begin{frame}[plain]
  \titlepage
\end{frame}

\begin{frame}{Outline}
  \tableofcontents
\end{frame}

%=======================================================================
\section{Motivation and data}
%=======================================================================
\begin{frame}{Why temporary help services?}
  \begin{itemize}
    \item Temporary help agencies (\textbf{NAICS 561320}) place workers on their own
          payroll at client firms. The sector is a bellwether: firms add and shed
          temps before permanent staff.
    \item It is also a large, low-wage segment of the labour market and the single
          biggest piece of \emph{Administrative and Support Services} (NAICS 56).
    \item Questions addressed in this analysis:
      \begin{enumerate}
        \item How has national temp employment moved across the 2008--09 recession,
              the 2010s expansion, and the pandemic?
        \item How large is temp help \emph{relative to} its parent sector?
        \item How do temp wages compare with the sector and with all private
              employment --- and is the gap changing?
        \item Who works in the sector (sex, age)?
      \end{enumerate}
  \end{itemize}
\end{frame}

\begin{frame}{Data and method}
  \begin{columns}[T]
    \begin{column}{0.54\textwidth}
      \textbf{Source.} Census \emph{Quarterly Workforce Indicators} (QWI), LEHD
      program, via the API endpoint \texttt{timeseries/qwi/sa}. Quarterly.
      QWI publishes no seasonally-adjusted series (confirmed against both the
      API and the raw bulk files); employment-level charts here use a
      self-computed approximation instead (see \emph{Seasonality}).

      \smallskip
      \textbf{Geography.} State-level counts for the 50 states + DC, summed to a
      national total. Earnings are \emph{employment-weighted} across states.

      \smallskip
      \textbf{Coverage.} 2005--2024 where available; the temp-help series is
      complete from 2005. Suppressed cells (flags \texttt{N/D/S}) are dropped.

      \smallskip
      \textbf{Also.} Race / ethnicity from QWI's \texttt{rh} dataset; the
      occupational mix from BLS \emph{OEWS} (May 2025).
    \end{column}
    \begin{column}{0.44\textwidth}
      \begin{block}{Industries compared}
        \footnotesize
        \begin{tabular}{@{}ll@{}}
          \toprule
          561320 & Temporary Help Services \\
          56     & Admin.\ \& Support Services \\
          00     & Total private sector \\
          \midrule
          561311 & Employment Placement \\
          561312 & Executive Search \\
          561330 & Prof.\ Employer Orgs (PEO) \\
          \bottomrule
        \end{tabular}
      \end{block}
      \begin{block}{Key measures}
        \footnotesize
        \textbf{Emp} --- beginning-of-quarter employment\\
        \textbf{EarnS} --- avg.\ monthly earnings, full-quarter workers\\
        \textbf{EarnHirAS} --- avg.\ monthly earnings, new hires
      \end{block}
    \end{column}
  \end{columns}
\end{frame}

%=======================================================================
\section{National employment}
%=======================================================================
\begin{frame}{National temp employment, 2005--2023}
  \fig[height=0.72\textheight]{national_timeline.png}
  \source{QWI NAICS 561320, sum of 50 states + DC. Seasonally adjusted
  (self-computed --- QWI has no official SA series; see \emph{Seasonality}).}
\end{frame}

\begin{frame}{National temp employment --- what the series shows}
  \begin{itemize}
    \item \textbf{Great Recession:} fell from \num{2.53} million (2006\,Q4) to
          \num{1.78} million (2009\,Q2), a drop of about \textbf{30\%} --- far
          deeper than the total-employment decline.
    \item \textbf{2010s expansion:} steady recovery; the prior peak was surpassed
          by 2014 and employment reached about \textbf{\num{3.11} million}
          (2018\,Q3).
    \item \textbf{Pandemic:} \num{2.99} million (2020\,Q1) $\rightarrow$
          \num{2.34} million (2020\,Q3), $-22\%$ in two quarters, then a rapid
          rebound to an all-time high of \textbf{\num{3.14} million} (2022\,Q2).
    \item \textbf{2023 slowdown:} down about \textbf{16\%} from that peak to
          \num{2.63} million by 2023\,Q4 --- back to 2013 levels, and consistent
          with temp help leading the broader labour-market cycle.
    \item Seasonally adjusted throughout (see \emph{Seasonality}, below) ---
          the raw series has a strong Q4 high / Q2 low every year.
  \end{itemize}
\end{frame}

%=======================================================================
\section{Seasonality}
%=======================================================================
\begin{frame}{Why seasonally adjusted --- and how}
  \fig[height=0.74\textheight]{national_seasonality.png}
\end{frame}

\begin{frame}{Seasonality --- how we approximate it}
  \begin{itemize}
    \item QWI has no official seasonally-adjusted series --- we checked both
          the API (\texttt{seasonadj} returns data only for \texttt{U},
          Unadjusted) and the raw bulk release files (same result).
    \item Every employment-level chart in this deck instead uses a
          self-computed approximation: classical multiplicative seasonal
          decomposition (\texttt{statsmodels.seasonal\_decompose}, period 4)
          fit to each series, dividing out the repeating quarterly pattern.
    \item The pattern is consistent and large: national temp employment runs
          \textbf{+3.1\%} above trend in Q4 (holiday-quarter staffing) and
          \textbf{$-$2.8\%} below trend in Q2.
    \item Ratios (temp's \% share of NAICS 56, turnover rates) are left on
          the raw QWI values --- seasonality in the numerator and denominator
          largely cancels there.
    \item This is an approximation, not an official series; treat exact
          turning-point quarters and small quarter-to-quarter moves with
          appropriate caution.
  \end{itemize}
\end{frame}

%=======================================================================
\section{Share of the parent sector}
%=======================================================================
\begin{frame}{Temp help as a share of Administrative \& Support Services (NAICS 56)}
  \fig[height=0.72\textheight]{sector_share_plot.png}
  \source{QWI: NAICS 561320 vs.\ NAICS 56 employment, national, seasonally
  adjusted (self-computed). The share (\%) table opposite uses the raw ratio.}
\end{frame}

\begin{frame}{Share of NAICS 56 --- roughly one-third, and cyclical}
  \begin{columns}[T]
    \begin{column}{0.46\textwidth}
      \small
      \begin{tabular}{@{}lS[table-format=2.1]@{}}
        \toprule
        Year & {Temp share of NAICS 56 (\%)} \\
        \midrule
        2010 & 28.7 \\
        2013 & 32.1 \\
        2015 & 33.8 \\
        2019 & 32.7 \\
        2020 & 30.4 \\
        2022 & 33.1 \\
        2023 & 30.7 \\
        \bottomrule
      \end{tabular}
    \end{column}
    \begin{column}{0.52\textwidth}
      \begin{itemize}
        \item Temp help is consistently \textbf{28--34\%} of all Administrative
              \& Support Services employment.
        \item The share \emph{rose} through the early 2010s recovery, then
              flattened near one-third.
        \item It falls in downturns (2020, 2023): temp jobs are cut first, so
              the sub-sector shrinks faster than its parent.
      \end{itemize}
    \end{column}
  \end{columns}
\end{frame}

%=======================================================================
\section{Geographic variation}
%=======================================================================
\begin{frame}{Where temp work concentrates}
  \fig[width=\textwidth,height=0.82\textheight,keepaspectratio]{temp_penetration_map.png}
\end{frame}

\begin{frame}{Penetration by state --- a Southeastern / industrial-belt story}
  \begin{columns}[T]
    \begin{column}{0.46\textwidth}
      \small
      \begin{tabular}{@{}lS[table-format=1.1]@{}}
        \toprule
        \multicolumn{2}{@{}l}{\textbf{Highest, 2023}} \\
        \midrule
        Tennessee      & 3.3 \\
        South Carolina & 3.2 \\
        Illinois       & 3.1 \\
        Alabama        & 3.0 \\
        Georgia        & 2.8 \\
        \midrule
        \multicolumn{2}{@{}l}{\textbf{Lowest, 2023}} \\
        \midrule
        South Dakota & 0.7 \\
        Hawaii       & 0.9 \\
        Vermont      & 1.1 \\
        \bottomrule
      \end{tabular}
    \end{column}
    \begin{column}{0.52\textwidth}
      \begin{itemize}
        \item Temp help is \textbf{2.2\%} of private employment nationally, but
              ranges roughly \textbf{5-fold} across 2023 states --- 0.7\% to 3.3\%.
        \item The high-penetration cluster is the Southeast and lower Midwest
              (TN, SC, AL, GA, NC, KY, IN, OH) --- auto, appliance, food-processing
              and warehousing supply chains that staff heavily through agencies.
        \item Penetration rose in \textbf{most states} from 2010 to 2023
              (e.g. MS $+0.7$, NE $+1.0$ pp).
        \item MI (shown at 2021) and AK (2016) have since stopped reporting to
              QWI; AK's 2016 penetration, \textbf{0.5\%}, was already the
              lowest in the country.
      \end{itemize}
    \end{column}
  \end{columns}
\end{frame}

%=======================================================================
\section{Wages}
%=======================================================================
\begin{frame}{Earnings: temp agencies vs.\ sector vs.\ all private}
  \fig[height=0.72\textheight]{wage_trends.png}
  \source{QWI EarnS / EarnHirAS, employment-weighted national averages, 2010--2023.}
\end{frame}

\begin{frame}{The temp wage gap --- large, but narrowing}
  \fig[height=0.60\textheight]{wage_gap.png}
  \vspace{-0.5em}
  \source{Temp EarnS as a percentage of the comparison group; dotted lines are period means.}
\end{frame}

\begin{frame}{Wages --- key numbers}
  \begin{columns}[T]
    \begin{column}{0.52\textwidth}
      \small
      \textbf{Average monthly earnings (EarnS)}
      \begin{tabular}{@{}lS[table-format=4.0]S[table-format=4.0]S[table-format=4.0]@{}}
        \toprule
        Year & {Temp} & {NAICS 56} & {Private} \\
        \midrule
        2010 & 2274 & 2934 & 3960 \\
        2015 & 2594 & 3311 & 4550 \\
        2019 & 3022 & 3797 & 5087 \\
        2022 & 4261 & 4783 & 6071 \\
        2023 & 4269 & 4926 & 6250 \\
        \bottomrule
      \end{tabular}
    \end{column}
    \begin{column}{0.46\textwidth}
      \small
      \textbf{Temp earnings as \% of\dots}
      \begin{tabular}{@{}lS[table-format=2.0]S[table-format=2.0]@{}}
        \toprule
        Year & {NAICS 56} & {Private} \\
        \midrule
        2010 & 78 & 57 \\
        2019 & 80 & 59 \\
        2023 & 87 & 68 \\
        \bottomrule
      \end{tabular}
      \medskip
      \begin{itemize}\footnotesize
        \item Period averages (2010--23): \textbf{81\%} of the sector,
              \textbf{60\%} of all private.
        \item The gap narrowed most in \textbf{2020--2022}, when a tight labour
              market bid up temp pay.
        \item For \emph{new hires}, temp pay is much closer to parity ---
              \(\approx\)97\% of sector, 89\% of private in 2023.
      \end{itemize}
    \end{column}
  \end{columns}
\end{frame}

%=======================================================================
\section{Turnover and job stability}
%=======================================================================
\begin{frame}{How transient is temp work?}
  \fig[height=0.74\textheight]{turnover_trends.png}
\end{frame}

\begin{frame}{Turnover --- temp workers churn about 3.5$\times$ the economy-wide rate}
  \begin{columns}[T]
    \begin{column}{0.54\textwidth}
      \fig[width=\textwidth]{turnover_comparison.png}
    \end{column}
    \begin{column}{0.44\textwidth}
      \small
      2023, per quarter (\%):
      \begin{tabular}{@{}lS[table-format=2.0]S[table-format=2.0]S[table-format=2.0]@{}}
        \toprule
                             & {Temp} & {N.56} & {Priv.} \\
        \midrule
        Stable emp.\ share   & 64 & 79 & 88 \\
        Churn rate           & 64 & 36 & 18 \\
        Stable-hire share    & 26 & 35 & 49 \\
        \bottomrule
      \end{tabular}
      \medskip
      \begin{itemize}\footnotesize
        \item Agencies make hires equal to \(\approx\)\textbf{2.5$\times$} their
              headcount each year; only \textbf{a quarter} of those hires last a
              full quarter.
        \item Churn \emph{drops} in recessions (2009, 2020): the loosely-attached
              are let go first, so who remains is more stable.
        \item The stable-hire share \emph{fell} through the 2010s --- temp jobs
              became a bit more spot-market, not less.
      \end{itemize}
    \end{column}
  \end{columns}
\end{frame}

%=======================================================================
\section{Demographics}
%=======================================================================
\begin{frame}{Who works in temp help? --- Sex (2023\,Q1, national)}
  \fig[width=\textwidth]{temp_employment_by_sex_2023_Q1.png}
  \bigskip
  \begin{itemize}
    \item Male \textbf{49.2\%} / Female \textbf{50.8\%} --- essentially even.
    \item Close to the overall private-sector split, unlike the male-dominated
          image of light-industrial staffing.
  \end{itemize}
  \source{QWI NAICS 561320, sex breakdown, 49 states with unsuppressed data.}
\end{frame}

\begin{frame}{Who works in temp help? --- Age (2023\,Q1, national)}
  \begin{columns}[c]
    \begin{column}{0.62\textwidth}
      \fig[width=\textwidth]{temp_employment_by_agegrp_2023_Q1.png}
    \end{column}
    \begin{column}{0.36\textwidth}
      \small
      \begin{tabular}{@{}lS[table-format=2.1]@{}}
        \toprule
        Age & {\%} \\
        \midrule
        14--24 & 13.3 \\
        25--34 & 27.0 \\
        35--44 & 22.3 \\
        45--54 & 17.1 \\
        55+    & 20.3 \\
        \bottomrule
      \end{tabular}
      \medskip
      \begin{itemize}\footnotesize
        \item About \textbf{two-thirds} are prime-age 25--54.
        \item Younger than the workforce overall, but not a youth sector ---
              one in five workers is 55+.
      \end{itemize}
    \end{column}
  \end{columns}
\end{frame}

\begin{frame}{Who works in temp help? --- Race and ethnicity (2023\,Q1, national)}
  \begin{columns}[c]
    \begin{column}{0.58\textwidth}
      \fig[width=\textwidth]{temp_employment_by_race_2023_Q1.png}
    \end{column}
    \begin{column}{0.40\textwidth}
      \small
      \begin{tabular}{@{}lS[table-format=2.1]@{}}
        \toprule
        Group & {\% of temp help} \\
        \midrule
        White               & 62.6 \\
        Black                & 26.7 \\
        Asian                & \phantom{0}6.2 \\
        Other / two+         & \phantom{0}4.5 \\
        \midrule
        Hispanic (any race)  & 21.1 \\
        \bottomrule
      \end{tabular}
      \medskip
      \begin{itemize}\footnotesize
        \item \textbf{Black workers are \(\approx\)27\%} of temp help --- about
              \textbf{twice} their \(\approx\)13\% share of US employment overall.
        \item \textbf{Hispanic} workers are \(\approx\)21\%, a smaller over-index.
        \item Consistent with the sector's concentration in warehousing,
              food-processing and light manufacturing.
      \end{itemize}
    \end{column}
  \end{columns}
  \source{QWI \texttt{rh} dataset, NAICS 561320, 49 reporting states. Race and Hispanic ethnicity are tabulated separately.}
\end{frame}

%=======================================================================
\section{Occupational mix}
%=======================================================================
\begin{frame}{What temp-agency workers actually do}
  \fig[width=\textwidth,height=0.80\textheight,keepaspectratio]{oews_occupation_mix.png}
\end{frame}

\begin{frame}{Occupational mix --- warehouse and factory labour, then everything else}
  \begin{columns}[T]
    \begin{column}{0.56\textwidth}
      \fig[width=\textwidth]{oews_top_occupations.png}
    \end{column}
    \begin{column}{0.42\textwidth}
      \begin{itemize}
        \item \textbf{Half} of temp-help jobs are blue-collar: transportation
              \& material moving (24\%) plus production (17\%) alone are
              \textbf{41\%}.
        \item One occupation --- \textbf{hand laborers / freight \& material
              movers} --- is \textbf{15\%} of the entire sector (394k jobs).
        \item But the sector is genuinely mixed: registered nurses (85k),
              substitute teachers (77k), HR specialists (75k), software
              developers (26k).
        \item The two largest groups are also the \textbf{lowest-paid}
              ($\approx\$37$k mean vs.\ the \$54k all-occupation mean) --- the
              mix, not just the "temp" label, drags sector pay down.
      \end{itemize}
    \end{column}
  \end{columns}
  \source{BLS OEWS, NAICS 561320, May 2025. OEWS annual wage assumes full-year work.}
\end{frame}

%=======================================================================
\section{Within employment services}
%=======================================================================
\begin{frame}{Employment-services sub-sectors (California, 2005--2024)}
  \fig[height=0.70\textheight]{employment_services_over_time.png}
  \source{QWI NAICS 561311/561320/561312/561330, California, seasonally adjusted (self-computed).}
\end{frame}

\begin{frame}{Employment services --- temp help dominates}
  \begin{itemize}
    \item In California, \textbf{561320 (temp help)} runs roughly
          \textbf{3--5$\times$} the size of the rest of employment services
          combined --- about \num{310}k jobs in 2005, an all-time peak of
          \num{410}k (2022\,Q2), then a fall to \num{325}k by 2024.
    \item \textbf{561330 (PEOs)} grew from \(\approx\)\num{48}k to \(\approx\)\num{58}k,
          with a step-up after 2021.
    \item \textbf{561311 (employment placement)} is flat near \num{44}k;
          \textbf{561312 (executive search)} is tiny (\(\approx\)\num{3}k).
    \item The single-state script (\texttt{temp\_agencies\_over\_time.py}) tracks
          the same California 561320 series, mirroring the national pattern:
          $-38\%$ in 2008--09, the same 2022\,Q2 all-time peak, the 2020 dip,
          and a slide from \num{410}k (2022) toward \num{325}k (2024).
  \end{itemize}
\end{frame}

%=======================================================================
\section{Takeaways}
%=======================================================================
\begin{frame}{Takeaways}
  \small
  \begin{enumerate}
    \item \textbf{Amplified cycle.} Temp employment swings far more than total
          employment --- $-30\%$ in 2008--09, $-22\%$ in early 2020, and a
          $16\%$ decline from the 2022 peak through 2023 that led the broader
          slowdown.
    \item \textbf{Stable structural weight.} Despite the swings, temp help holds
          near \textbf{one-third} of Administrative \& Support Services.
    \item \textbf{Persistent but shrinking wage gap.} Temp workers earn about
          \textbf{80\%} of sector pay and \textbf{60\%} of all-private pay on
          average; both ratios rose markedly in 2020--2022 and remained higher
          through 2023. New-hire pay is close to parity.
    \item \textbf{Very short attachment.} Quarterly churn near \textbf{65\%}
          (vs.\ 18\% economy-wide) and only a quarter of temp hires last a full
          quarter --- the defining feature of the sector.
    \item \textbf{Concentrated geographically.} A \textbf{5-fold} spread across
          states, highest in the Southeast / industrial belt where supply chains
          staff through agencies.
    \item \textbf{Blue-collar core.} Half the jobs are warehouse / material-moving
          and production work --- and those two groups are also the lowest-paid,
          which is much of why the sector's wages sit low.
    \item \textbf{Distinctive workforce.} Roughly even by sex and prime-aged,
          but \textbf{Black workers are about twice as prevalent} as in the
          economy overall (27\% vs.\ 13\%); Hispanic workers modestly so.
  \end{enumerate}
\end{frame}

\begin{frame}{Caveats}
  \begin{itemize}
    \item QWI is a by-product of state UI records; small state $\times$ quarter
          $\times$ demographic cells are suppressed, so demographic totals rest
          on 48--49 states, not all 51 jurisdictions.
    \item Employment-level charts are seasonally adjusted with a
          self-computed approximation, not an official series (see
          \emph{Seasonality}); wage, turnover, and share ratios remain on
          raw QWI values, where seasonality in the numerator and denominator
          largely cancels.
    \item ``Employment'' is a count of jobs (beginning-of-quarter), not
          full-time-equivalents; temp assignments are short, so headcount
          overstates labour input relative to other sectors.
    \item Education is published only on the QWI \texttt{se} dataset for the
          25+ restriction, and every state-level cell for NAICS 561320 is
          suppressed --- so no education breakdown is possible.
    \item \textbf{Michigan} left QWI after 2021 and \textbf{Alaska} after 2016;
          the state map falls back to each state's last available year and
          tags the tile accordingly.
    \item OEWS is a separate survey (establishment-reported, annual) and is not
          in the BLS API --- the 561320 table was taken from the OEWS query
          system. Its "annual wage" assumes full-year work.
    \item California is used for the sub-sector and single-state views; national
          patterns may differ.
  \end{itemize}
\end{frame}

\begin{frame}[plain]
  \begin{center}
    \Large Thank you\\[1em]
    \normalsize Code and figures: \texttt{code/} and \texttt{output/} in the
    \texttt{Temp\_agencies} repository.
  \end{center}
\end{frame}

\end{document}
